\section{The original line over an arbitrary relative affine base}

\begin{lemma}[A mixed projective expression for the original line]
\label{lem:kn-relative-arithmetic-original-line-expression}
\lean{EnvelopeProperty.Vanishing.RelativeArithmeticLineExpression.exists_original_line_expression}
\uses{thm:kn-arithmetic-original-line-generated}
\leanok
Let $k$ be a field, let $R$ be any commutative ring, let $X$ be a
$k$-scheme, and let $j:X\hookrightarrow\mathbb P_R^d$ be a closed
immersion.  For every invertible module $L$ on $X$, there are integers
$N,b\geq0$ and a morphism $h:X\to\mathbb P_k^b$ over $k$ with
\[
 h^*\mathcal O_{\mathbb P_k^b}(1)
   \simeq L\otimes j^*\mathcal O_{\mathbb P_R^d}(N).
\]
Consequently the same original $L$ is isomorphic to
$h^*\mathcal O(1)\otimes(j^*\mathcal O(N))^\vee$.
Neither $R$ nor $X$ is assumed integral, and $X$ need not be proper
over $k$.
\end{lemma}
\begin{proof}
The original line $L$ is quasicoherent and of finite type.  Serre
generation for the closed immersion $j$ gives an $N$ such that
$L\otimes j^*\mathcal O(N)$ has finitely many global generating
sections.  This tensor product is again invertible.  Projectivize
those same finitely many sections using the field structure of $X$.
On the open where one chosen section generates the line, divide all
the sections by that generator.  The resulting coefficients define
a morphism to the corresponding standard affine chart of
$\mathbb P_k^b$.  On pairwise intersections their coordinate ratios
agree, so these maps glue to a morphism $h$ over $k$.  The same local
generator identifies the pulled-back $\mathcal O(1)$ with the
generated line.  These identifications agree on intersections since
they carry each coordinate section to its original global section;
hence they glue to the displayed whole-line isomorphism.  Tensor
with the dual of $j^*\mathcal O(N)$, use associativity, and cancel the
line with its dual by evaluation to obtain the final assertion.
\end{proof}
